% mnras_template.tex 
%
% LaTeX template for creating an MNRAS paper
%
% v3.0 released 14 May 2015
% (version numbers match those of mnras.cls)
%
% Copyright (C) Royal Astronomical Society 2015
% Authors:
% Keith T. Smith (Royal Astronomical Society)

% Change log
%
% v3.2 July 2023
%	Updated guidance on use of amssymb package
% v3.0 May 2015
%    Renamed to match the new package name
%    Version number matches mnras.cls
%    A few minor tweaks to wording
% v1.0 September 2013
%    Beta testing only - never publicly released
%    First version: a simple (ish) template for creating an MNRAS paper

%%%%%%%%%%%%%%%%%%%%%%%%%%%%%%%%%%%%%%%%%%%%%%%%%%
% Basic setup. Most papers should leave these options alone.
\documentclass[fleqn,usenatbib]{mnras}

% MNRAS is set in Times font. If you don't have this installed (most LaTeX
% installations will be fine) or prefer the old Computer Modern fonts, comment
% out the following line
\usepackage{newtxtext,newtxmath}
\usepackage{multirow}
\usepackage[round-mode=places,
            separate-uncertainty=true,
            table-align-uncertainty=true]{siunitx}
\usepackage[table]{xcolor}
% Depending on your LaTeX fonts installation, you might get better results with one of these:
%\usepackage{mathptmx}
%\usepackage{txfonts}

% Use vector fonts, so it zooms properly in on-screen viewing software
% Don't change these lines unless you know what you are doing
\usepackage[T1]{fontenc}

% Allow "Thomas van Noord" and "Simon de Laguarde" and alike to be sorted by "N" and "L" etc. in the bibliography.
% Write the name in the bibliography as "\VAN{Noord}{Van}{van} Noord, Thomas"
\DeclareRobustCommand{\VAN}[3]{#2}
\let\VANthebibliography\thebibliography
\def\thebibliography{\DeclareRobustCommand{\VAN}[3]{##3}\VANthebibliography}


%%%%% AUTHORS - PLACE YOUR OWN PACKAGES HERE %%%%%

% Only include extra packages if you really need them. Avoid using amssymb if newtxmath is enabled, as these packages can cause conflicts. newtxmatch covers the same math symbols while producing a consistent Times New Roman font. Common packages are:
\usepackage{graphicx}	% Including figure files
\usepackage{amsmath}	% Advanced maths commands

%%%%%%%%%%%%%%%%%%%%%%%%%%%%%%%%%%%%%%%%%%%%%%%%%%

%%%%% AUTHORS - PLACE YOUR OWN COMMANDS HERE %%%%%

% Please keep new commands to a minimum, and use \newcommand not \def to avoid
% overwriting existing commands. Example:
%\newcommand{\pcm}{\,cm$^{-2}$}	% per cm-squared

%%%%%%%%%%%%%%%%%%%%%%%%%%%%%%%%%%%%%%%%%%%%%%%%%%

%%%%%%%%%%%%%%%%%%% TITLE PAGE %%%%%%%%%%%%%%%%%%%

% Title of the paper, and the short title which is used in the headers.
% Keep the title short and informative.
\title[Correction to: The cosmic dipole in Quaia]{
    Correction to: The cosmic dipole in the Quaia sample of quasars:
    a Bayesian analysis}

% The list of authors, and the short list which is used in the headers.
% If you need two or more lines of authors, add an extra line using \newauthor
\author[Mittal, Oayda $\&$ Lewis]{
Vasudev Mittal,$^{1,2*}$\thanks{E-mail: vasudeviiser@gmail.com}
Oliver T. Oayda,$^{2*\thanks{E-mail: oliver.oayda@sydney.edu.au}}$
and Geraint F. Lewis${^2}$
\\
% List of institutions
$^{*}$Joint first author\\
$^{1}$Department of Physical Sciences, IISER Mohali, Knowledge City, Sector 81, SAS Nagar, Manauli PO 140306, Punjab, India\\
$^{2}$Sydney Institute for Astronomy, School of Physics A28, The University of Sydney, NSW 2006, Australia\\
}

% These dates will be filled out by the publisher
\date{Accepted XXX. Received YYY; in original form ZZZ}

% Enter the current year, for the copyright statements etc.
\pubyear{2024}

% Don't change these lines
\begin{document}
\label{firstpage}
\pagerange{\pageref{firstpage}--\pageref{lastpage}}
\maketitle

% Select between one and six entries from the list of approved keywords.
% Don't make up new ones.
\begin{keywords}
    errata, addenda -- quasars: general -- cosmology: observations
    -- cosmology: theory -- large-scale structure of Universe
\end{keywords}

%%%%%%%%%%%%%%%%%%%%%%%%%%%%%%%%%%%%%%%%%%%%%%%%%%

%%%%%%%%%%%%%%%%% BODY OF PAPER %%%%%%%%%%%%%%%%%%

\noindent In the original paper \citep{mittal2024},
 an error was made in the calculation of the distribution
of spectral indices $\alpha$,
as appearing in fig. 2 of that paper.
The error arises from equation~(3),
which in the original paper reads
\begin{equation}
    m_\nu = -2.5 \log_{10} S_\nu + \text{ZP}. \label{eq:mag}
\end{equation}
While this equation is mathematically correct,
the issue relates to the correct interpretation of the units
of the passband flux density $S_\nu$ when using \textit{Gaia} magnitudes.
These magnitudes, like the G band magnitude as
appearing in the \verb|phot_g_mean_mag| column
of the \textit{Gaia} DR3 release table,
are defined such that $S_\nu$ in Equation~\eqref{eq:mag} is in \emph{units of photoelectrons per second}
\citep[$e^- s^{-1}$;][]{hambly2022}.
We noted this in our original paper at the beginning
of Section~4.2.2,
at which stage we mentioned that to convert from $e^- s^{-1}$ into Jy,
one must first multiply the passband flux density by
a correction factor $c_\nu$ to give units of W\,m$^{-2}$\,Hz$^{-1}$.
However, since this factor was not applied in equation~(4) in the original work,
equation~(5) is incorrect as one cannot use the relation
$S_\nu \propto \nu^{-\alpha}$ if $S_\nu$ is in units of $e^- s^{-1}$.

To correct this,
let $S'_\nu$ denote the passband flux density in
W\,m$^{-2}$\,Hz$^{-1}$ such that $S'_\nu = c_\nu S_\nu$.
Substituting this into Equation~\eqref{eq:mag} yields
\begin{equation}
m_\nu = -2.5 \log_{10} S'_\nu + 2.5 \log_{10} c_\nu + \text{ZP}.
\end{equation}
Thus, equations~(4) and (5) in the original work
should be replaced with these corrected versions:
\begin{align}
    m_{\mathrm{G - BP}} &= 2.5 \log_{10}
    \left( S'_{\mathrm{BP}} / S'_{\mathrm{G}} \right) +
    2.5 \log_{10}
        \left( c_{\mathrm{G}} / c_{\mathrm{BP}} \right) + k \\
    % m_{\mathrm{G - BP}} &= 2.5 \left[ \log_{10} (c_{\mathrm{BP}} S_{\mathrm{BP}}) - \log_{10}( c_{\mathrm{G}} S_{\mathrm{G}}) \right] + k \\
    \implies \alpha_i &= \frac{\epsilon + k - m_{\mathrm{G - BP}}}{2.5 \log_{10} (\nu_{\mathrm{BP}} / \nu_{\mathrm{G}})} \label{eq:spectral-index}
\end{align}
where $\epsilon$ is $2.5 \log_{10} (c_{\mathrm{G}} / c_{\mathrm{BP}} )$ and $k$ is $\text{ZP}_{\mathrm{G}} - \text{ZP}_{\mathrm{BP}}$ as defined originally.
The value of these correction factors for each passband
are given in section~5.4.1 of the \textit{Gaia} DR3 documentation \citep{busso2022}.
With the correction applied,
fig. 2 in the original work is to be replaced with Fig.~\ref{fig:spectral_indices} as shown below.
Since fig. 3 in the original work
included the correction factor when converting \textit{Gaia} magnitudes
to flux densities,
this figure does not need to be altered.
\begin{figure}
    \centering
    \includegraphics[width=\linewidth]{Figures/fig2_corrected.pdf}
    \caption{Distribution of spectral indices $\alpha$ in the Quaia low and Quaia high samples computed from $m_{G - BP}$. Blue and orange are used to denote Quaia low and Quaia high respectively. The mean spectral indices $\bar{\alpha}$ for each catalogue are indicated by the vertical lines.}
    \label{fig:spectral_indices}
\end{figure}

As the \citet{ellis1984} relation is dependent on $\alpha$,
the correction to the spectral index
changes the expected dipole amplitude $\mathcal{D}$.
We repeated the methodology in Section~4.2.2 with the amended
values of $\alpha$,
finding a distribution of dipole amplitudes as shown in Fig.~\ref{fig:dipole-amplitude}.
\begin{figure}
    \centering
    \includegraphics[width=\linewidth]{Figures/fig4_corrected.pdf}
    \caption{Probability distribution for the dipole amplitude assuming $v_\text{CMB}$. This follows from the analysis in Section~4.2.2 and was performed with both the Quaia low and Quaia high samples.}
    \label{fig:dipole-amplitude}
\end{figure}
This is to replace fig.~4 of the original work.
We thus take the mean amplitude of $\mathcal{\bar{D}} = 0.0048$ for Quaia low
and $\mathcal{\bar{D}} = 0.0043$ for Quaia high,
using these as the dipole expectations.

Correcting for this error in the way described above impacts the relative Bayesian evidences
of our different hypotheses.
This is because models which incorporate the expected dipole amplitude will
necessarily arrive at different marginal likelihoods.
These are models $M_5$ (kinematic velocity) and $M_6$ (kinematic dipole)
in our original paper.
Accordingly,
the last two rows of tables A1--A4 in the original paper
are to be replaced with those given here in
Tables~\ref{tab:bayes-low-p2p}--\ref{tab:bayes-high-poisson}.

Based on these tables, our conclusions are altered slightly.
The relative ordering of each model in terms of evidential power is not changed.
Namely,
for the Quaia low sample with a $40^\circ$ galactic plane mask,
the kinematic dipole (model $M_6$) is still the prevailing model.
However, the difference in support between $M_6$ and the next-favoured model,
$M_4$ (the marginal likelihood of which is given in the original work),
becomes $\ln B_{64} = 1.5$.
This means that the kinematic dipole is somewhat less preeminent than as in the original work
($\ln B_{65} = 2.6$).

Additionally, the conclusion in section~6.3 in the original work
needs to be replaced.
We now find that our conclusions are, to a certain extent,
sensitive to the choice of prior.
As mentioned there,
after adjusting the prior function for $\mathcal{D}$
from $\mathcal{U}[0,1.0]$ to $\mathcal{U}[0,0.1]$,
the marginal likelihood of model $M_4$ increases to 8.0
(Quaia low; $b^\circ < |40|$ masked).
Based on the corrected Bayes factors in Table~\ref{tab:bayes-low-p2p},
this means that model $M_4$ (kinematic direction) is the prevailing model
with $\ln B_{46} = 0.7$,
in which case we find $\mathcal{D} \approx 11\substack{+5 \\ -5} \times 10^{-3}$.
Although, this Bayes factor is indicative of only a slight preference
for $M_4$ over $M_6$.
Meanwhile,
the conclusions for Quaia high using the adjusted prior on $\mathcal{D}$
are unchanged.

Though these aspects have been altered,
it still remains that the Quaia sample
strongly favours a dipole aligning with the direction
of the CMB dipole.
The question of its amplitude is not as decisive;
while we cannot rule out the possibility of a value
larger than expected from the CMB dipole,
we have also pointed towards support for
it being equal to the CMB value.
This leaves significant scope for future inquiry
to attempt to resolve this matter.

% the conclusion that we are heading towards the kinematic expectation
% is the same
% the orientation of the dipole is aligned; how strongly can one
% say that the amplitude has kinematic origin i.e. corresponds
% to the value derived from the CMB.?

\begin{table*}
  \centering
   \begin{tabular}{l  S[table-format=7.1, round-precision=1] S[table-format=7.1,round-precision=1] S[table-format=7.1, round-precision=1] S[table-format=7.1,round-precision=1] S[table-format=7.1, round-precision=1] S[table-format=7.1, round-precision=1] }
\hline
\multirow{2}*{Hypothesis} & \multicolumn{6}{c}{Galactic mask angle $b^\circ$}\\
  & 0 & 10 & 20 &30 & 40 & 30*\\
  \hline
  $M_{5}$ (Kinematic Velocity) & 30.4 & 27.8 & 17.7 & 9.1  & 5.3  & 7.5  \\
  $M_{6}$ (Kinematic Dipole)   & 16.2 & 18.4 & 15.7 & 11.0 & 7.3\cellcolor{black!10} & 10.1 \\

\hline
 \end{tabular}
 \caption{Table of Bayes Factors for different hypotheses and galactic masks using the Quaia low catalogue with the point-by-point analysis. Here, 30$^*$ represents the combination of a $30^\circ$ mask and a $4\,\text{sr}$ circular mask centered at the $(l^\circ, b^\circ)=(0,0)$. The highlighted cell represents the model with the highest Bayes factor, indicating it has the strongest level of support.}
 \label{tab:bayes-low-p2p}
 \end{table*}

 \begin{table*}
  \centering
   \begin{tabular}
   {l  S[table-format=7.1, round-precision=1] S[table-format=7.1,round-precision=1] S[table-format=7.1, round-precision=1] S[table-format=7.1,round-precision=1] S[table-format=7.1, round-precision=1] S[table-format=7.1, round-precision=1] }
\hline
\multirow{2}*{Hypothesis} & \multicolumn{6}{c}{Galactic mask angle $b^\circ$}\\
  & 0 & 10 &20&30&40&30*\\
  \hline
$M_{5}$ (Kinematic Velocity) & 30.4 & 27.9 & 17.7 & 9.0  & 5.1 & 7.4 \\
$M_{6}$ (Kinematic Dipole)   & 16.3 & 18.4 & 15.6 & 11.0 & 7.1\cellcolor{black!10} & 10.1 \\
\hline
 \end{tabular}
 \caption{As for Table~\ref{tab:bayes-low-p2p} but with the Poisson statistics.}
 \label{tab:bayes-low-poisson}
 \end{table*}

 \begin{table*}
  \centering
   \begin{tabular}
   {l  S[table-format=7.1, round-precision=1] S[table-format=7.1,round-precision=1] S[table-format=7.1, round-precision=1] S[table-format=7.1,round-precision=1] S[table-format=7.1, round-precision=1] S[table-format=7.1, round-precision=1] }
\hline
\multirow{2}*{Hypothesis} & \multicolumn{6}{c}{Galactic mask angle $b^\circ$}\\
  & 0 & 10 &20&30&40&30*\\
  \hline
$M_{5}$ (Kinematic Velocity) & 62.9 & 56.6 & 37.3 & 20.9 & 14.6 & 11.5 \\
$M_{6}$ (Kinematic Dipole)   & 21.8 & 25.2 & 24.5 & 18.5 & 14.5 & 12.6 \\
\hline
 \end{tabular}
 \caption{As for Table~\ref{tab:bayes-low-p2p} but with Quaia high.}
 \label{tab:bayes-high-p2p}
 \end{table*}

 \begin{table*}
  \centering
   \begin{tabular}
   {l  S[table-format=7.1, round-precision=1] S[table-format=7.1,round-precision=1] S[table-format=7.1, round-precision=1] S[table-format=7.1,round-precision=1] S[table-format=7.1, round-precision=1] S[table-format=7.1, round-precision=1] }
\hline
\multirow{2}*{Hypothesis} & \multicolumn{6}{c}{Galactic mask angle $b^\circ$}\\
  & 0 & 10 &20&30&40&30*\\
  \hline
$M_{5}$ (Kinematic Velocity) & 63.1 & 56.5 & 37.4 & 20.7 & 14.9 & 11.3 \\
$M_{6}$ (Kinematic Dipole)   & 21.9 & 25.1 & 24.2 & 18.6 & 14.5 & 12.7 \\
\hline
 \end{tabular}
  \caption{As for Table~\ref{tab:bayes-low-p2p} but with Quaia high and the Poisson statistics.}
  \label{tab:bayes-high-poisson}
 \end{table*}


\section*{Acknowledgements}
We extend our sincere gratitude to Nathan Secrest,
who pointed out the steepness of our original spectral indices
and generously shared some of their own results.
This led us to uncover the error mentioned in this correction.

This work has made use of data from the European Space Agency (ESA) mission \textit{Gaia} (\url{https://www.cosmos.esa.int/gaia}),
processed by the \textit{Gaia} Data Processing and Analysis Consortium (DPAC, \url{https://www.cosmos.esa.int/web/gaia/dpac/consortium}).
Funding for the DPAC has been provided by national institutions,
in particular the institutions participating in the \textit{Gaia} Multilateral Agreement.

%%%%%%%%%%%%%%%%%%%% REFERENCES %%%%%%%%%%%%%%%%%%

% The best way to enter references is to use BibTeX:

\bibliographystyle{mnras}
\bibliography{quaia_correction} % if your bibtex file is called example.bib

%%%%%%%%%%%%%%%%%%%%%%%%%%%%%%%%%%%%%%%%%%%%%%%%%%


% Don't change these lines
\bsp	% typesetting comment
\label{lastpage}
\end{document}

% End of mnras_template.tex
